\section{Results of the historical baselines}
\label{sec:baselines}
\subsection{Sanity references and centralized learning}
Phase 0 established random-ranking and raw-popularity references under the
same full-catalog evaluator later used for learned models. The historical
nonprivate raw-popularity test NDCG@10 is $0.044292$; random ranking is much
lower, approximately $0.0038$. These references help distinguish personalized
learning from a pipeline that merely returns popular items or fails to learn.
The full exact sanity metrics are retained in the curated results archive.

% Generated by scripts/generate_tables.py; do not edit numeric cells.
\begin{table}[tbp]\centering\small
\caption{Matched centralized and federated controls (three seeds).}\label{tab:central}
\setlength{\tabcolsep}{4pt}
\begin{tabular}{lrrrr}\toprule
Model & NDCG@10 & Seed SD & HR@10 & MRR@10 \\
\midrule
B0 & 0.085972 & 0.000456 & 0.171975 & 0.060051 \\
B0-UW & 0.091745 & 0.001866 & 0.181882 & 0.064536 \\
B1 & 0.087539 & 0.002074 & 0.164190 & 0.064527 \\
\bottomrule\end{tabular}
\tabnote{Seeds 42/123/2026. This B1 mean differs from the five-seed B1 mean used in later comparisons. Source: centralized\_federated\_3seeds.csv.}
\end{table}

\Cref{tab:central} gives the final three-seed comparison. Centralized B0 reaches
$0.085972$, B0-UW $0.091745$, and B1 $0.087539$ test NDCG@10. The paired
B0-UW-minus-B0 difference is $0.005773$, with ordinary 95\% interval
$[-0.000164,0.011910]$. B1-minus-B0-UW is $-0.004206$, with interval
$[-0.010903,0.002288]$. The total B1-minus-B0 difference is $0.001567$,
with interval $[-0.005353,0.008668]$. All three overall intervals contain zero.

The point estimates suggest that user weighting partly offsets the cost of
federated optimization, but this study does not resolve an overall superiority
claim for either change. In particular, B1's numerical proximity to B0 should
not be used to assert that federation has no effect: the intervention changes
weighting, optimizer, and training organization together. B0-UW improves the
interpretability of that comparison while preserving the uncertainty.

Historical personalization checks replace a user's own embedding with another
user's embedding and substantially reduce utility. The saved B1 seed-42 check
falls from approximately $0.0890$ to $0.0202$ NDCG@10. This is a sanity check
that the model uses user-specific information, not a five-seed estimate of a
new serving method. E1 later broadens this diagnostic on validation states.

\subsection{Nonprivate low rank}
% Generated by scripts/generate_tables.py; do not edit numeric cells.
\begin{table}[tbp]\centering\small
\caption{Converged nonprivate federated utility (five seeds).}\label{tab:nonprivate}
\setlength{\tabcolsep}{4pt}
\begin{tabular}{lrr}\toprule
Model & NDCG@10 $\pm$ seed SD & Adjusted interval vs. B1 \\
\midrule
B1 & $0.08901\pm0.00249$ & Reference \\
B3 r4 & $0.05947\pm0.00381$ & [-0.041779, -0.017893] \\
B3 r8 & $0.07101\pm0.00596$ & [-0.027511, -0.008949] \\
B3 r16 & $0.07633\pm0.00722$ & [-0.020950, -0.004560] \\
B3 r32 & $0.06807\pm0.01354$ & [-0.031093, -0.011160] \\
\bottomrule\end{tabular}
\tabnote{Intervals are 98.75\% paired-user bootstrap intervals for the four-comparison B3 family. Source: final\_5seed\_means.csv and paired\_contrasts.csv.}
\end{table}

The five-seed converged B1 mean is $0.089008$, which differs from the three-seed
B1 entry because the seed set differs. The B3 means at ranks 4, 8, 16, and 32
are $0.059468$, $0.071010$, $0.076333$, and $0.068074$. All four adjusted
intervals favor B1 (\cref{tab:nonprivate}). Low rank therefore incurs a utility
cost before DP is introduced under these retained settings.

The rank trend is not monotonic. Rank 16 has the largest B3 mean; rank 32 has
larger seed variation and documented convergence/stability complications.
This is evidence about the trained models, not a theorem that rank 32 has
less representational capacity. Optimization, selected rates, and stopping
behavior contribute to the outcomes. The expanded rank-32 diagnostic ceiling
is disclosed rather than hidden inside an apparently common training budget.

\reportfigure{01_nonprivate_controls.pdf}{Nonprivate controls and low-rank utility}
{Left: matched three-seed centralized and federated test NDCG@10. Right:
five-seed converged B1 and B3. Whiskers show seed SD. The panels have different
seed sets. Overall weighting-decomposition intervals include zero; all four
adjusted B3-versus-B1 intervals are negative. Sources: retained centralized
control and final five-seed tables.}{fig:nonprivate}

\subsection{Private full matrix and low rank}
% Generated by scripts/generate_tables.py; do not edit numeric cells.
\begin{table}[tbp]\centering\small
\caption{Historical private and matched no-DP test NDCG@10 (five-seed means).}\label{tab:historical}
\setlength{\tabcolsep}{4pt}
\begin{tabular}{lrrrrr}\toprule
Target / control & B2 Full & B4 r4 & B4 r8 & B4 r16 & B4 r32 \\
\midrule
No DP & 0.057889 & 0.040834 & 0.048083 & 0.049143 & 0.050931 \\
8 & 0.052271 & 0.038544 & 0.045385 & 0.043628 & 0.044022 \\
4 & 0.042169 & 0.037846 & 0.042653 & 0.038274 & 0.044489 \\
2 & 0.025428 & 0.029079 & 0.034485 & 0.033993 & 0.036504 \\
1 & 0.011461 & 0.015976 & 0.018162 & 0.016685 & 0.015390 \\
\bottomrule\end{tabular}
\tabnote{All rows use 50 rounds. Historical tuning opportunities differ. Seed SDs and secondary metrics appear in the supplement. Source: final\_5seed\_means.csv.}
\end{table}

The full-matrix B2 mean falls from $0.057889$ in its matched 50-round no-DP
control to $0.052271$, $0.042169$, $0.025428$, and $0.011461$ at privacy
targets 8, 4, 2, and 1. Its no-DP value is already below converged B1, so
the entire B1--B2 gap cannot be assigned to privacy noise. B4's no-DP controls
also start below B2's no-DP value.

The sixteen adjusted B4-minus-B2 comparisons give a conditional crossover:
all four ranks lose at target eight; no adjusted difference is detected at
target four; ranks 8, 16, and 32 win at target two; and rank 8 wins at target
one. These are the four positive comparisons. Several additional ordinary
95\% intervals look favorable but do not survive the declared adjustment.
The validation-selected rank at target one is 16, not the test-best rank 8;
the former's adjusted interval includes zero. Selecting the latter after
examining test scores would change the interpretation of the experiment.

At target two, for example, B4 rank 32 reaches $0.036504$ against B2's
$0.025428$. Yet both are below the historical raw-popularity point reference.
The low-rank model's smaller loss from its own weaker no-DP control does not
by itself establish better recommendation quality. This observation motivates
the user-level private popularity comparator introduced in E1.

\reportfigure{02_historical_private_utility.pdf}{Historical private utility}
{Five-seed test NDCG@10 with seed-SD whiskers for B2 and all B4 ranks at 50
rounds. Larger epsilon means weaker privacy. The dotted popularity line is
nonprivate raw popularity. Historical tuning budgets differ. Four of sixteen
adjusted B4-minus-B2 comparisons are positive; line segments only connect
discrete experimental settings.}{fig:historical}

\subsection{Communication and the original hypothesis}
The rank-8 two-factor payload is 111,744 bytes per participating client-round,
compared with 861,184 bytes for Full. At a fixed 50-round horizon, this
approximately $7.7$-fold payload reduction also produces a corresponding
reduction in modeled total tensor volume. B2 averages $4.026$ GB and B4 rank
8 about $0.522$ GB. These are meaningful savings in the simulator's accounting.

Converged training has a different comparison. B1 averages $93.911$ GB to its
best checkpoint; B3 rank 32 averages $180.329$ GB despite a smaller per-round
payload. The latter requires many more rounds. Rank 4 and rank 8 do reduce
total volume, but their utility is lower. The full communication table is in
\cref{app:algorithms} and the supplement. No experiment in this project
optimizes utility under an identical end-to-end communication cap.

The historical results therefore partially support the initial hypothesis:
low rank can offer a better private point on the utility/payload trade-off,
especially when the full model is heavily degraded. They reject its strongest
form, namely that fewer shared noisy parameters reliably improve utility.
Unequal historical tuning and multiple changes in optimization prevent a
dimension-only causal interpretation. E1 was designed to examine the proposed
noise explanation more directly.
\FloatBarrier
